\documentclass[11pt]{article}
\usepackage[margin=1in]{geometry}
\usepackage[T1]{fontenc}
\usepackage{lmodern}
\usepackage{amsmath,amssymb}
\usepackage{hyperref}

\title{Literature-Implied Status for arXiv:1310.6874}
\author{Run fa\_banach\_001, agent\_lane\_16}
\date{July 18, 2026}

\begin{document}
\maketitle

\paragraph{Status.}
This note records a literature-implied partial answer to an open problem stated in Daniel Kornlein's
\emph{Quantitative results for Halpern iterations of nonexpansive mappings}, arXiv:1310.6874.
It is not a new result of the run. The arbitrary uniformly smooth case remains outside the scope of the supporting theorem.

\paragraph{Source problem.}
Kornlein obtains a rate of metastability for Halpern iterations in uniformly smooth Banach spaces relative to a rate of metastability for the resolvent path $(z_t)$.
Near the end of the introduction, after the display describing the form of the Halpern bound, the paper states that it is still an open problem to extract a rate of metastability for $(z_t)$ in the uniformly smooth case, and notes that such a result would include the spaces $L_p$, $1<p<\infty$, $p\neq 2$.

\paragraph{Supporting theorem and identification.}
Kohlenbach and Sipos, \emph{The finitary content of sunny nonexpansive retractions}, arXiv:1812.04940, extract a rate of metastability for the Reich approximating/resolvent curve in Banach spaces which are both uniformly convex and uniformly smooth. Their introduction states that this covers all $L^p$ spaces for $1<p<\infty$ and gives the first explicit rate of metastability for the Shioji--Takahashi extension of Wittmann's Halpern theorem to this class of spaces.

The paper's Section 7 then applies the extracted resolvent bound to Halpern iterations: it states that the earlier Kohlenbach--Leustean proof-mining analysis was modulo resolvent convergence, and that the new paper is in a position to complete that analysis under the additional hypothesis that $X$ is uniformly convex. The displayed theorem in that section gives, for uniformly convex and uniformly smooth $X$, a bound $N\leq\Sigma(\varepsilon,\omega_\tau,g,b,\Theta_{b,\eta,\tau,\theta,\alpha,\gamma},\beta_1,\beta_2,\beta_3)$ such that $\|x_m-x_n\|\leq\varepsilon$ for all $m,n\in[N,N+g(N)]$.

\paragraph{Scope.}
This answers the source problem for the important uniformly convex-and-uniformly smooth subcase, in particular for all $L_p$, $1<p<\infty$. It does not settle the full uniformly smooth case without uniform convexity.

\paragraph{Provenance.}
This is a literature-implied identification. The supporting paper explicitly cites Kornlein's arXiv:1310.6874 in its bibliography and explicitly presents a Halpern-iteration consequence, but the relation to the exact source sentence is made here by matching Kornlein's missing resolvent-metastability input with the Kohlenbach--Sipos resolvent-metastability theorem.

\paragraph{References.}
\begin{itemize}
\item D. Kornlein, \emph{Quantitative results for Halpern iterations of nonexpansive mappings}, arXiv:1310.6874; J. Math. Anal. Appl. 428 (2015), no. 2, 1161--1172.
\item U. Kohlenbach and A. Sipos, \emph{The finitary content of sunny nonexpansive retractions}, arXiv:1812.04940; Comm. Contemp. Math. 23 (2021), no. 3, 1950093.
\item U. Kohlenbach and L. Leustean, \emph{On the computational content of convergence proofs via Banach limits}, Phil. Trans. R. Soc. A 370 (2012), 3449--3463.
\end{itemize}

\end{document}
